% Section 2.3, from lectures/L02/README.md: the objectives, the questions,
% the further reading, and what the next chapter does with all of it.


\section{Review}
\label{c2:sec:review}

\needspace{6\baselineskip}
\subsection*{What you should be able to do now}
After this chapter, you should be able to:
\begin{itemize}
  \item Compute a time constant, and say how many of them a given settling accuracy needs.
  \item State what a phasor is, and what has to be true of a circuit before one is allowed.
  \item Write down the impedance of a capacitor and an inductor at any frequency, from memory.
  \item Sketch a first-order response from its corner frequency alone, magnitude and phase.
  \item Say what the phase is doing a decade either side of the corner, to within a degree or two.
  \item Convert between a ratio and decibels without a calculator for the cases that matter.
  \item Extend a real solver from real to complex arithmetic without duplicating it.
\end{itemize}

\needspace{6\baselineskip}
\subsection*{Questions to test yourself}
You should be able to answer:
\begin{enumerate}
  \item A capacitor's impedance falls with frequency and an inductor's rises. At what frequency are they equal, and what is special about a circuit operating there?
  \item Why is $j\omega$ a legitimate substitute for $d/dt$, and what property of the circuit does that argument depend on?
  \item A first-order low-pass is 3 dB down at its corner. How far down is it half a decade above?
  \item At what frequency has a first-order low-pass reached 45 degrees of phase lag, and how much magnitude has it lost there?
  \item A colleague says a filter ``does nothing'' a decade below its corner. What has it done to the phase there, and when would that matter?
  \item Your AC solver returns the right magnitude and a phase of the wrong sign at every frequency. Name the most likely single cause.
  \item An amplifier settles to within 1 per cent in 5 microseconds. How long does it need for 0.01 per cent, assuming one time constant dominates?
\end{enumerate}

\needspace{6\baselineskip}
\subsection*{Further reading}
\begin{itemize}
  \item \emph{The Art of Electronics}, Horowitz and Hill, chapter 1, sections on reactance and RC circuits, for a second treatment of the same material.
\end{itemize}

\needspace{6\baselineskip}
\subsection*{Looking ahead}
\begin{itemize}
  \item Chapter~\ref{c3:ch:filters} puts the reactance to work. Passive filters, what cascading two of them actually gives you, and then the operational amplifier as a black box with two rules.
\end{itemize}
